% eksperymentalne tłumaczenie na włoski

%

\subsection{Obroty} % lang-pl
\subsection{Rotazioni} % lang-it

\begin{definition}
[obrót] % lang-pl
[rotazione] % lang-it
    Niech $A$ będzie punktem, zaś $\alpha$ miarą kąta. % lang-pl
    Funkcję $R_A^\alpha$ zadaną następująco: $R_A^\alpha(X) = Y$ wtedy i tylko wtedy, gdy kąt $\angle XAY$ ma miarę $\alpha$ i $|AY| = |AX|$ nazywamy obrotem o kąt $\alpha$ wokół punktu $A$ zwanego środkiem obrotu. % lang-pl
    Siano $A$ un punto e $\alpha$ un angolo. % lang-it
    Chiamiamo rotazione di angolo $\alpha$ attorno al punto $A$, detto centro di rotazione, la funzione $R_A^\alpha$ definita come segue: $R_A^\alpha(X) = Y$ se e solo se l'angolo $\angle XAY$ ha ampiezza $\alpha$ e $|AY| = |AX|$. % lang-it
\end{definition}

Oczywiście i tutaj można podać wzór. % lang-pl
Wyznacza się go najpierw dla obrotów wokół zera, by następnie zauważyć, że dowolny obrót sprowadza się do tamtego przy użyciu dwóch translacji (jednej, by przenieść się na początek i drugiej, by wrócić): % lang-pl
Naturalmente anche in questo caso si può fornire una formula esplicita. % lang-it
La si ricava dapprima per le rotazioni attorno all'origine e poi si osserva che ogni rotazione si riconduce a quel caso mediante due traslazioni (una per portarsi all'origine e una per tornare indietro): % lang-it

\begin{equation}
    R_A^\alpha(x,y) = \begin{pmatrix}
        x_A + (x_0 - x_A) \cos \alpha - (y_0 - y_A) \sin \alpha\\
        y_A + (x_0 - x_A) \sin \alpha - (y_0 - y_A) \cos \alpha
    \end{pmatrix}.
\end{equation}

Symetrie środkowe nazywa się (rzadko) półobrotem, czyli obrotem o kąt $\pi$; robi tak Eves \cite[s. 105]{eves1_1972}. % lang-pl
Le simmetrie centrali vengono talvolta chiamate semirotazioni, cioè rotazioni di angolo $\pi$; così fa Eves \cite[p. 105]{eves1_1972}. % lang-it

\begin{proposition}
    Złożenie dwóch symetrii osiowych, których osie przecinają się pod kątem $\alpha$ jest obrotem o kąt $2 \alpha$ wokół punktu przecięcia wspomnianych osi. % lang-pl
    La composizione di due simmetrie assiali i cui assi si intersecano formando un angolo $\alpha$ è una rotazione di angolo $2\alpha$ attorno al punto d'intersezione di tali assi. % lang-it
\end{proposition}

To uogólnienie faktu \ref{two_axial_one_point}, tam było $\alpha = \pi$, a półobrót (obrót o kąt $\pi$) jest tym samym, co symetria środkowa. % lang-pl
Si tratta di una generalizzazione del fatto \ref{two_axial_one_point}: lì si aveva $\alpha = \pi$ e una semirotazione (rotazione di angolo $\pi$) coincide con una simmetria centrale. % lang-it

\begin{proposition}
\label{for_banach_11}%
    Niech $A \neq B$ będą dwoma punktami, zaś $\alpha, \beta$ dwiema miarami kątów. % lang-pl
    Jeśli $\alpha + \beta$ nie jest wielokrotnością kąta pełnego, $2 \pi$, to $R_B^\beta \circ R_A^\alpha = R_C^{\alpha + \beta}$, gdzie $C$ jest takim punktem, by trójkąt $\triangle ABC$ miał kąty $\alpha/2$ przy $A$, $\beta / 2$ przy $B$. % lang-pl
    Siano $A \neq B$ due punti e $\alpha, \beta$ due angoli. % lang-it
    Se $\alpha + \beta$ non è un multiplo dell'angolo giro $2\pi$, allora $R_B^\beta \circ R_A^\alpha = R_C^{\alpha + \beta}$, dove $C$ è tale che il triangolo $\triangle ABC$ abbia angolo $\alpha/2$ in $A$ e angolo $\beta/2$ in $B$. % lang-it

    W przeciwnym przypadku złożenie tych dwóch obrotów jest translacją. % lang-pl
    In caso contrario, la composizione di queste due rotazioni è una traslazione. % lang-it
\end{proposition}

%